\documentclass[a4paper, 11pt]{ctexart}

% ==========================================
% 导言区：常用宏包配置
% ==========================================
\usepackage[margin=2.5cm]{geometry} % 页面边距
\usepackage{amsmath, amsfonts, amssymb} % 数学公式支持
\usepackage{graphicx} % 插入图片
\usepackage{float} % 图片表格浮动体控制
\usepackage{hyperref} % 交叉引用与超链接
\hypersetup{
    colorlinks=true,
    linkcolor=blue,
    filecolor=magenta,      
    urlcolor=cyan,
    citecolor=black
}
\usepackage{booktabs} % 三线表
\usepackage{listings} % 插入代码
\usepackage{xcolor} % 代码高亮颜色

% 代码格式设置
\lstset{
    basicstyle=\ttfamily\small,
    keywordstyle=\color{blue},
    commentstyle=\color{green!50!black},
    stringstyle=\color{red},
    breaklines=true,
    frame=single,
    numbers=left,
    numberstyle=\tiny\color{gray}
}

% ==========================================
% 标题信息
% ==========================================
\title{基于子空间低秩重建的2D磁共振指纹成像仿真与验证}
\author{BME1309.01 生医工影像导论期末项目 \\ \\ \textit{刘恒希、秦子丰、肖智健}}
\date{\today}

\begin{document}

\maketitle

% 摘要（可选，但技术报告通常建议有一段简短的执行摘要）
\begin{abstract}
% 摘要撰写要求：用3-4句话概括项目的核心方法（TGAS轨迹、子空间LLR重建）、主要工作（仿真跑通+真实数据验证）以及最终结论。尽量精简。
\end{abstract}

\newpage
\tableofcontents
\newpage

% ==========================================
% 正文开始
% ==========================================

\section{项目背景与出发点}
% 内容说明：
% - 传统MRI的局限性（单一定性加权图像、耗时长）。
% - MRF（磁共振指纹技术）的优势（快速、多参数定量成像）。
% - 指出当前MRF面临的挑战（欠采样伪影、重建计算量大），引出本项目的切入点。

\section{项目目的}
% 内容说明：
% - 核心目标：在2D场景下，利用特定算法改进MRF图像质量，缩短采集时间。
% - 具体任务分解：轨迹设计、字典生成与压缩、低秩重建、定量图匹配。

\section{核心原理与理论依据}
% 内容说明：参考曹晓智等人的相关论文，简要罗列项目依赖的物理和数学理论，切忌长篇大论，把公式和变量说明白即可。
\subsection{扩展相位图 (EPG) 与字典生成}
% - 说明如何模拟不同T1、T2在特定TR下的信号演化。
\subsection{TGAS 螺旋轨迹与 k 空间采样}
% - 解释可变密度螺旋轨迹及 Tiny-Golden-Angle-Shuffling 旋转采样的原理和优势。
\subsection{子空间与局部低秩 (LLR) 重建}
% - 解释奇异值分解(SVD)降维的过程。
% - 给出重建的目标函数公式（数据一致性项 + LLR正则化项），说明如何消除欠采样伪影。
\subsection{模板匹配 (Template Matching)}
% - 简述如何通过内积最大化从压缩字典中匹配出最终的 T1、T2、PD 参数。

\section{代码实现手段}
% 内容说明：直奔重点，描述代码架构和数据流向。
\subsection{开发环境与依赖}
% - 简要提一句基于 Python 的统一包结构，以及用到的核心库（SigPy, Torch 等）。
\subsection{模块化流水线设计}
% - 结合你们的 pipeline 流程，分点说明（可使用 bullet points）：
%   1. 轨迹与体模生成模块
%   2. 字典仿真与压缩模块
%   3. 前向模拟模块（加噪）
%   4. 重建与参数映射模块

\section{实验流程与数据说明}
% 内容说明：交代验证算法所用的数据来源。
\subsection{仿真数据准备}
% - 提及构建的 2D Phantom 及其组织参数设定。
\subsection{真实扫描数据来源}
% - 简短说明真实数据的来源（如开源的 cMRF 扫描数据）、数据结构（HDF5格式、FA文本）等。
\subsection{算法检验与评估方法}
% - 说明如何评判重建结果的好坏（例如：与 Reference 图像比对、观察伪影消除情况、对比不同迭代参数下的表现）。

\section{参数调试与结果分析}
% 内容说明：报告的核心部分。直接上图表和数据，用客观事实说话。
\subsection{模拟参数运行与验证}
% - 展示纯模拟流程下的跑通结果。
% - 展示不同关键参数（如：迭代次数 \texttt{n-iter}、LLR正则化权重 \texttt{lambda-llr}、子空间维度等）对重建图像质量的影响。
% - 简要分析哪组参数在去除伪影和保留细节上达到了最佳平衡。
\subsection{真实扫描数据测试}
% - 在模拟数据验证算法无误后，引入真实数据的表现。
% - 展示基于真实数据的重建定量图（T1/T2图）。
% - 客观分析真实数据带来的新挑战（如噪声分布更复杂、潜在的物理伪影等）以及 LLR 算法在真实数据上的鲁棒性。

\section{总结与展望}
% 内容说明：
% - 总结已完成的核心工作和达成的效果（LLR迭代发挥作用）。
% - 坦诚目前存在的不足（如：真实数据结果仍需优化、未加入 $B_0$ 校正等）。
% - 提出后续改进方向。

% ==========================================
% 参考文献
% ==========================================
\newpage
\begin{thebibliography}{99}
% 在此处添加具体的参考文献引用，如项目依据的论文
\bibitem{cao2022} 
Xiaozhi Cao, Congyu Liao, et al. \textit{Optimized multi-axis spiral projection MR fingerprinting with subspace reconstruction for rapid whole-brain high-isotropic-resolution quantitative imaging}. Magnetic Resonance in Medicine, 2022.
\end{thebibliography}

\end{document}